%% ch02_systems_design.tex -- Intelligence Engineering, chapter 02
%% "Introduction to Machine Learning Systems Design".
%% Spine transferred from Huyen, Designing Machine Learning Systems, chapter
%% 2. The subchapters and the frame titles under them are the book's own
%% section and subsection headings, in the book's order.
%%
%% STATE: titles and TEMPLATE layout only. No body text. Every frame carries
%% the layout it is meant to be filled into, and a % TEMPLATE comment saying
%% why that layout and what belongs in each slot. Filling them is the next pass.
%%
%% Fragment of intelligence_engineering.tex. One chapter is one lecture and
%% gets exactly ONE \lecturepage, at the top. Inside it every \subsection
%% opens on the subchapter divider, which the master emits automatically via
%% \AtBeginSubsection; do not write those divider frames by hand.

\begin{refsection}

\lecturepage{II}{Introduction to Machine Learning Systems Design}{}
\section[Systems Design]{Introduction to Machine Learning Systems Design}

% ----------------------------------------------------------------------
\subsection{Business and ML Objectives}

% TEMPLATE: two block contrast. The heading names both sides, so each side
% takes one box rather than one shared list, levelled by the equal height
% group. Left box takes the metrics the business actually reports, right box
% the metrics a model is trained to move. The filling pass has to make the
% mapping from the right box to the left one explicit.
\begin{frame}{Business and ML Objectives}
  \begin{columns}[T,onlytextwidth]
    \begin{column}{0.47\textwidth}
      \begin{tdblockt}[equal height group=obj]{Business objectives}
      \end{tdblockt}
    \end{column}
    \begin{column}{0.47\textwidth}
      \begin{tdblockt}[equal height group=obj]{ML objectives}
      \end{tdblockt}
    \end{column}
  \end{columns}
  \slidesources{Huyen2022}
\end{frame}

% ----------------------------------------------------------------------
\subsection{Requirements for ML Systems}

% TEMPLATE: withmargin plus tddef. This is the first of the book's four
% requirements and the book states it as a definition, so it takes the
% slide's one titled box. Margin holds the aside that makes the requirement
% hard for a machine learning system: it can fail silently and still return
% a prediction that looks plausible.
\begin{frame}{Reliability}
  \begin{withmargin}
    \begin{tdmain}
      \begin{tddef}{Reliability}
      \end{tddef}
    \end{tdmain}
    \begin{tdmargin}
    \end{tdmargin}
  \end{withmargin}
  \slidesources{Huyen2022}
\end{frame}

% TEMPLATE: three plain boxes in a row. Scalability is not one thing in the
% book, it is growth along several axes at once, so each axis takes one box
% and the row makes the count visible instead of burying it in a list. The
% boxes stay unheaded until the filling pass names the axes; the equal
% height group keeps the row level.
\begin{frame}{Scalability}
  \begin{columns}[T,onlytextwidth]
    \begin{column}{0.31\textwidth}
      \begin{tdblock}[equal height group=scale]
      \end{tdblock}
    \end{column}
    \begin{column}{0.31\textwidth}
      \begin{tdblock}[equal height group=scale]
      \end{tdblock}
    \end{column}
    \begin{column}{0.31\textwidth}
      \begin{tdblock}[equal height group=scale]
      \end{tdblock}
    \end{column}
  \end{columns}
  \slidesources{Huyen2022}
\end{frame}

% TEMPLATE: withmargin, and a deliberate break from the row above it. Main
% column takes what maintainability asks of the artifacts themselves, code,
% data and documentation. Margin takes who those artifacts are handed to,
% the book's point that an ML system is built by people who do not share a
% vocabulary.
\begin{frame}{Maintainability}
  \begin{withmargin}
    \begin{tdmain}
    \end{tdmain}
    \begin{tdmargin}
    \end{tdmargin}
  \end{withmargin}
  \slidesources{Huyen2022}
\end{frame}

% TEMPLATE: tddef alone. The book states this requirement as a definition
% too, and the frame reads stronger bare than boxed inside a margin layout,
% so the term carries the whole slide. The box takes the capacities the book
% asks for and nothing else; what they imply for the workflow is the next
% subchapter's job.
\begin{frame}{Adaptability}
  \begin{tddef}{Adaptability}
  \end{tddef}
  \slidesources{Huyen2022}
\end{frame}

% ----------------------------------------------------------------------
\subsection{Iterative Process}

% TEMPLATE: side image hero, the first of this chapter's two figures. The
% heading is a process, so the book's own cycle owns the frame and the grey
% placeholder holds its place until the diagram is drawn. sidebody takes the
% steps in the book's order, the phimgside caption the one line that says why the
% arrows run backwards as often as forwards.
\begin{frame}{Iterative Process}
  \phimgside[0.33]{iterative_process.png}{}
  \begin{sidebody}
  \end{sidebody}
  \slidesources{Huyen2022}
\end{frame}

% ----------------------------------------------------------------------
\subsection{Framing ML Problems}

% TEMPLATE: two by two grid of plain boxes, the opener for the four frames
% below it. The heading introduces a countable set, so the grid is sized to
% that count and each box previews one of the four headings that follow
% rather than arguing anything itself. Unheaded on purpose: the detail
% frames own the names.
\begin{frame}{Types of ML Tasks}
  \begin{columns}[T,onlytextwidth]
    \begin{column}{0.47\textwidth}
      \begin{tdblock}[equal height group=tasks]
      \end{tdblock}
      \vskip0.5em
      \begin{tdblock}[equal height group=tasks]
      \end{tdblock}
    \end{column}
    \begin{column}{0.47\textwidth}
      \begin{tdblock}[equal height group=tasks]
      \end{tdblock}
      \vskip0.5em
      \begin{tdblock}[equal height group=tasks]
      \end{tdblock}
    \end{column}
  \end{columns}
  \slidesources{Huyen2022}
\end{frame}

% TEMPLATE: two block contrast. The heading is a versus, so each task takes
% one box with its header read straight off the heading, levelled by the
% equal height group. Each box takes what that task predicts in the book's
% words, and the filling pass owes the book's reminder that either one can
% be turned into the other.
\begin{frame}{Classification versus regression}
  \begin{columns}[T,onlytextwidth]
    \begin{column}{0.47\textwidth}
      \begin{tdblockt}[equal height group=clsreg]{Classification}
      \end{tdblockt}
    \end{column}
    \begin{column}{0.47\textwidth}
      \begin{tdblockt}[equal height group=clsreg]{Regression}
      \end{tdblockt}
    \end{column}
  \end{columns}
  \slidesources{Huyen2022}
\end{frame}

% TEMPLATE: comparison table. A second pair of boxes in a row would stall,
% and the book separates these two along named dimensions rather than in
% prose, so the columns are the two sides from the heading and the empty
% rows are for those dimensions, one per row. booktabs only, no vertical
% rules.
\begin{frame}{Binary versus multiclass classification}
  \footnotesize
  \renewcommand{\arraystretch}{1.35}
  \begin{tabular}{@{}>{\raggedright\arraybackslash}p{0.24\textwidth} >{\raggedright\arraybackslash}p{0.32\textwidth} >{\raggedright\arraybackslash}p{0.32\textwidth}@{}}
    \toprule
     & \textbf{Binary} & \textbf{Multiclass}\\
    \midrule
     & & \\
     & & \\
     & & \\
     & & \\
    \bottomrule
  \end{tabular}
  \slidesources{Huyen2022}
\end{frame}

% TEMPLATE: two code panels. The book settles this distinction on how a
% label is encoded, so the contrast belongs in the mono panels where a label
% vector can be shown literally instead of described. The captions are the
% two sides from the heading; each panel takes the encoding the book prints
% for that case, and the ambiguity it creates.
\begin{frame}{Multiclass versus multilabel classification}
  \begin{columns}[T,onlytextwidth]
    \begin{column}{0.47\textwidth}
      \begin{tdcodet}{Multiclass}
      \end{tdcodet}
    \end{column}
    \begin{column}{0.47\textwidth}
      \begin{tdcodet}{Multilabel}
      \end{tdcodet}
    \end{column}
  \end{columns}
  \slidesources{Huyen2022}
\end{frame}

% TEMPLATE: two plain boxes. The heading promises more than one framing of
% the same problem, so the framings sit side by side where they can be read
% off against each other. Unheaded, because the heading names no framing;
% each box takes one framing with its inputs, its outputs and what happens
% to it when the set of classes grows.
\begin{frame}{Multiple ways to frame a problem}
  \begin{columns}[T,onlytextwidth]
    \begin{column}{0.47\textwidth}
      \begin{tdblock}[equal height group=ways]
      \end{tdblock}
    \end{column}
    \begin{column}{0.47\textwidth}
      \begin{tdblock}[equal height group=ways]
      \end{tdblock}
    \end{column}
  \end{columns}
  \slidesources{Huyen2022}
\end{frame}

% TEMPLATE: withmargin, the opener for the frame below it. Main column takes
% what an objective function is for and the losses the book names, margin
% the aside that the objective is a choice and not a given. Kept general on
% purpose: the detail, two objectives inside one loss, is the next frame.
\begin{frame}{Objective Functions}
  \begin{withmargin}
    \begin{tdmain}
    \end{tdmain}
    \begin{tdmargin}
    \end{tdmargin}
  \end{withmargin}
  \slidesources{Huyen2022}
\end{frame}

% TEMPLATE: side image hero, the chapter's second and last figure. The
% heading is an architecture decision, one model per objective with the
% predictions combined afterwards, which reads as a diagram and not as
% prose. sidebody takes the argument for decoupling in the book's order,
% the phimgside caption what the decoupling costs.
\begin{frame}{Decoupling objectives}
  \phimgside[0.33]{decoupling_objectives.png}{}
  \begin{sidebody}
  \end{sidebody}
  \slidesources{Huyen2022}
\end{frame}

% ----------------------------------------------------------------------
\subsection{Mind Versus Data}

% TEMPLATE: two block contrast. This is the chapter's loudest versus and the
% book gives it as a standing argument with named people on both sides, so
% each camp takes one box, headed from the heading and levelled by the equal
% height group. Each box takes that camp's position and the evidence it
% leans on.
\begin{frame}{Mind Versus Data}
  \begin{columns}[T,onlytextwidth]
    \begin{column}{0.47\textwidth}
      \begin{tdblockt}[equal height group=mind]{Mind}
      \end{tdblockt}
    \end{column}
    \begin{column}{0.47\textwidth}
      \begin{tdblockt}[equal height group=mind]{Data}
      \end{tdblockt}
    \end{column}
  \end{columns}
  \slidesources{Huyen2022}
\end{frame}

% ----------------------------------------------------------------------
\subsection{Summary}

% TEMPLATE: takeaway alone. A summary earns no scaffold of its own, so the
% chapter closes on one centred line with nothing competing against it. The
% line is the single sentence the chapter should leave behind, lifted from
% the book's own summary rather than written fresh.
\begin{frame}{Summary}
  \begin{takeaway}
  \end{takeaway}
  \slidesources{Huyen2022}
\end{frame}

% ----------------------------------------------------------------------
% This chapter's own references. The refsection opened above scopes the
% citation numbering to this chapter, so chapter_pdfs/<this>.pdf resolves
% its own [n] markers instead of pointing at a list it does not contain.
\begin{frame}{References}
  \printbibliography[heading=none]
\end{frame}
\end{refsection}
